\par\medskip
\noindent
5.5. Para expresar cómodamente los factores extremos de las
filtraciones (5.1.10) y (5.1.11), en el caso de reducción semiestable,
introduzcamos los esquemas en grupos constantes torcidos sobre $S$
\begin{equation}\tag{5.5.1}
 \underline M'_o=D(T_o),\qquad \underline M_o=D(T'_o),
\end{equation}
duales de los toros maximales $T_o,T'_o$ de las fibras especiales de
$A,A'$ (SGA 3, X 5.7). Puesto que $S$ es local completo, estos esquemas
en grupos pueden considerarse las fibras especiales de esquemas
constantes torcidos
\begin{equation}\tag{5.5.2}
 \underline M',\qquad \underline M\quad\hbox{sobre }S.
\end{equation}
Así, si el toro $T_o$ es escindido (por ejemplo, si $k$ es
separablemente cerrado), es decir, si $\underline M'_o$ es el grupo
constante definido por el grupo ordinario
\begin{equation}\tag{5.5.3}
 M'=\operatorname{Hom}_{k\text{-}\mathrm{gr}}(T_o,\mathbb G_{mk}),
\end{equation}
entonces $\underline M'$ es el grupo constante definido por $M'$:
\begin{equation*}
 \underline M'=M'_S.
\end{equation*}
Puede observarse además que, como $T_o$ y $T'_o$ son isógenos (2.2.6),
$T'_o$ es escindido siempre que $T_o$ lo sea. En este caso conviene
introducir también
\begin{equation}\tag{5.5.4}
 M=\operatorname{Hom}_{k\text{-}\mathrm{gr}}(T'_o,\mathbb G_{mk}),
\end{equation}
y entonces
\begin{equation*}
 \underline M=M_S.
\end{equation*}
Con esta notación, para todo entero $n$ existen evidentemente
isomorfismos
\begin{equation}\tag{5.5.5}
 \underline M'_n\simeq D(T_n),\qquad
 \underline M_n\simeq D(T'_n),
\end{equation}
donde $T'_n$ se define en términos de $A'$ como $T_n$ en términos de
$A$ (5.1.2). Se deducen, por tanto, isomorfismos canónicos (donde
$\mathbb Z_\ell(1)=T_\ell(\mathbb G_{mS})$)
\begin{equation}\tag{5.5.6}
 T_\ell(T_n)\simeq
 \check{\underline M}'_n\otimes_{\mathbb Z}\mathbb Z_\ell(1),
 \qquad
 T_\ell(T'_n)\simeq
 \check{\underline M}_n\otimes_{\mathbb Z}\mathbb Z_\ell(1),
\end{equation}
y, pasando al límite sobre $n$, isomorfismos canónicos
\begin{equation}\tag{5.5.7}
 \left\{
 \begin{aligned}
  T_\ell(A^o)^t=T_\ell(\widehat T)
   &\simeq \check{\underline M}'\otimes\mathbb Z_\ell(1),\\[1ex]
  T_\ell(A^{\prime o})^t=T_\ell(\widehat T')
   &\simeq \check{\underline M}\otimes\mathbb Z_\ell(1).
 \end{aligned}
 \right.
\end{equation}
En el caso de reducción semiestable, la dualidad, junto con el teorema
de ortogonalidad 5.2 y la relación (5.2.2), proporciona entonces
isomorfismos canónicos
\begin{equation}\tag{5.5.8}
 \left\{
 \begin{aligned}
  T_\ell(A_K)/T_\ell(A_K)^f
   &\simeq \underline M_K\otimes_{\mathbb Z}\mathbb Z_\ell,\\[1ex]
  T_\ell(A'_K)/T_\ell(A'_K)^f
   &\simeq \underline M'_K\otimes_{\mathbb Z}\mathbb Z_\ell.
 \end{aligned}
 \right.
\end{equation}
Análogamente, la misma fórmula (5.2.2) proporciona isomorfismos
canónicos
\begin{equation}\tag{5.5.9}
 \left\{
 \begin{aligned}
  T_\ell(A_K)/T_\ell(A_K)^t
   &\simeq D\bigl(T_\ell(A'_K)^f\bigr),\\[1ex]
  T_\ell(A'_K)/T_\ell(A'_K)^t
   &\simeq D\bigl(T_\ell(A_K)^f\bigr),
 \end{aligned}
 \right.
\end{equation}
donde debe recordarse que los grupos $T_\ell(A_K)^f$ y
$T_\ell(A'_K)^f$ que aparecen a la derecha son las restricciones a
$\eta$ de los pro-$\ell$-grupos de Barsotti--Tate $T_\ell(A^o)^f$ y
$T_\ell(A^{\prime o})^f$, respectivamente, definidos sobre $S$. Aquí
$D$ denota la dualidad en el sentido de los progrupos de Barsotti--Tate,
a saber, $\underline{\operatorname{Hom}}(-,\mathbb Z_\ell(1))$.
Tenemos, pues:

\medskip
\noindent
\underline{Proposición} 5.6. \underline{Supongamos que $S$ es completo
y que $A_K$ posee reducción semiestable sobre $S$. Entonces, para todo
número primo $\ell$} (\underline{que puede ser igual a la característica residual
de} $S$), \underline{el pro}-$\ell$-\underline{grupo de Barsotti--Tate $T_\ell(A_K)^t$ es la
fibra genérica de un progrupo de Barsotti--Tate sobre $S$ «de tipo
tórico», dado por (5.5.7), mientras que
$T_\ell(A_K)/T_\ell(A_K)^f$ es la fibra genérica de un progrupo de
Barsotti--Tate sobre $S$ de tipo proétale, dado por (5.5.8). Por último,
los grupos $T_\ell(A_K)^f$ y $T_\ell(A_K)/T_\ell(A_K)^t$ son las fibras
genéricas de progrupos de Barsotti--Tate definidos sobre $S$, iguales,
respectivamente, a $T_\ell(A^o)^f$ y a
$D(T_\ell(A^{\prime o})^f)=
\underline{\operatorname{Hom}}(T_\ell(A^{\prime o})^f,
\mathbb Z_\ell(1))$.}

\medskip
\noindent
5.6.1. Obsérvese que, para un recubrimiento étale $S'$ de $S$, escindir
$T_\ell(A_K)^t$ equivale a escindir
$T_\ell(A_K)/T_\ell(A_K)^f$: estas condiciones expresan,
respectivamente, que $S'$ escinde los grupos constantes torcidos
$\sisi{\underline M}{\underline M'}$ y $\sisi{\underline M'}{\underline M}$, que son isógenos porque $T_o$ y
$T'_o$ lo son. Al mismo tiempo se ve que existe una menor extensión
étale $S'$ de $S$ que cumple este cometido.

Registramos también la siguiente consecuencia fácil de las definiciones
de 2.2.3.

